\thispagestyle{abstractpagestyle}

\vspace*{36pt}

\begin{center}
    \textbf{\fontsize{20pt}{24pt} \selectfont ABSTRACT}
\end{center}

\vspace{24pt}

Quantum image processing could speed up tasks such as filtering and edge detection, but today's noisy intermediate-scale quantum (NISQ) devices allow only shallow circuits on a handful of qubits, and they remain noticeably noisy. The flexible representation of quantum images (FRQI) stores each pixel intensity as a rotation angle on a single color qubit, selected by multi-controlled gates that read a register of address qubits. Preparing that state naively, with multi-controlled $X$ gates and no helper qubits, produces deep circuits with many controlled-NOT (CX) gates after compilation.

This thesis asks a narrow engineering question: does changing only \emph{how} the multi-controlled gates are built change the cost and the noise robustness of the same FRQI state? It compares two synthesis modes, a \emph{no-ancilla} decomposition (naive) and a reusable \emph{v-chain} pattern that borrows a few ancilla qubits, on $4\times 4$, $8\times 8$, and $16\times 16$ grayscale test images in the Qiskit Aer simulator. The downstream stage, the quantum Hadamard edge detector (QHED) scored against a classical Sobel reference, is held fixed so that any difference comes from the preparation circuit alone. Three kinds of measurement are reported side by side rather than in isolation. The resource side records transpiled qubits, depth, and CX count. Reconstruction quality under synthetic depolarizing noise is measured with PSNR, SSIM, and state fidelity, and edge agreement is the SSIM between the QHED map and the Sobel reference. A separate proof-of-concept corrects readout bias on the color qubit by inverting an empirical $2\times 2$ confusion matrix. Every number traces back to versioned CSV/JSON manifests and Python scripts in the project repository.

All experiments ran on simulators only, with no IBM hardware runs. A few of the $16\times 16$ v-chain results are missing because the test machine ran out of RAM.

\vfill
